\chapter{土壤卫生  【重点章节】}
\setcounter{chapter}{7}
\chapterintro{
本章为\keyword{重点考察章节}（3学时）。掌握土壤环境背景值的概念及意义；掌握土壤污染的基本特点；掌握土壤重金属污染（镉、铊、铬等）和农药污染的健康危害；掌握持久性有机污染物（POPs）的特性及危害；了解土壤卫生防护措施。
}

\section{土壤环境特征}

\subsection{土壤的组成}
\begin{itemize}[leftmargin=*]
    \item \imp{土壤固相}：矿物质（占土壤固体总质量90\%以上）+ 有机物质（腐殖质、生物残体、土壤生物）
    \item \imp{土壤液相}：水分及其水溶物质
    \item \imp{土壤气相}：土壤空隙中多种气体混合物；上层与大气相似，深层O$_2$减少、CO$_2$增加
\end{itemize}

\subsection{土壤物理学特征}

\begin{defbox}
\textbf{土壤质地}：按土壤粒级及其组合比例而定的土壤名称，反映土壤固相颗粒分布情况。分为\imp{砂土、壤土、黏土}。壤土通气透水、保水保肥性能均较好，是良好的土壤质地。

\textbf{土壤孔隙度（Soil Porosity）}：自然状态下，单位容积土壤中孔隙容积所占百分率。对土壤性质的影响：①容水性（保存水分的能力）；②渗水性（水分渗透的能力）；③透气性（空气穿透的能力）；④毛细管现象（水分沿孔隙上升的程度）。

\textbf{规律：}土壤颗粒越小，容水性越大，透气性越差。
\end{defbox}

\begin{keybox}
\textbf{土壤环境背景值（Background Level）}：该地区未受或少受人为污染影响的天然土壤中各种元素的含量。

\textbf{实践意义：}
\begin{enumerate}[leftmargin=*]
    \item 评价化学污染物对土壤污染程度的\imp{参照值}
    \item 制定土壤中有害化学物质卫生标准的\imp{重要依据}
    \item 评价土壤化学环境对居民健康影响的\imp{重要依据}
    \item 土地资源开发利用和\imp{地方病防治}的科学依据
\end{enumerate}
\end{keybox}

\begin{defbox}
\textbf{土壤环境容量/土壤负载容量}：一定土壤环境单元在一定时限内遵循环境质量标准，即维持土壤生态系统的正常结构与功能，保证农产品的生物学产量与质量，在不使环境系统污染的前提下，土壤环境所能容纳污染物的最大负荷量。是充分利用土壤环境的纳污能力、实现污染物总量控制、合理制定环境质量标准和卫生标准的重要依据。

\textbf{腐殖质（Humus）}：动植物残体经分解转化形成的有机物质，占土壤有机质总量的\imp{85\%--90\%}。
\end{defbox}

\begin{tipbox}
\textbf{土壤生物性污染监测指标：}
\begin{itemize}[leftmargin=*]
    \item \imp{大肠菌值}：发现大肠菌的最少土壤克数，代表人畜粪便污染的主要指标，也是代表肠道传染病危险性的主要指标
    \item \imp{产气荚膜杆菌值}：代表粪便污染的指标，与大肠菌在土壤中数量的消长关系可判定污染时间长短——产气荚膜杆菌多而大肠菌相对少→\textbf{陈旧性污染}；大肠菌值更多→\textbf{新鲜污染}（危害性较大）
    \item \imp{蛔虫卵数}：直接说明在流行病学上是否对人体健康有威胁。根据蛔虫卵在土壤中的不同发育阶段及活卵所占百分比，可判断土壤的自净程度
\end{itemize}
\end{tipbox}

\section{土壤的污染、自净及转归}

\subsection{土壤污染的基本特点}

\begin{keybox}
\textbf{四大基本特点：}
\begin{enumerate}[leftmargin=*]
    \item \imp{隐蔽性}——不易被直接察觉，往往通过食物链传递才被发现
    \item \imp{累积性与地域性}——污染物长期累积，区域差异大
    \item \imp{不可逆转性}——重金属等一旦进入极难清除
    \item \imp{治理的长期性}——修复周期长、成本高
\end{enumerate}
\end{keybox}

\subsection{污染方式}

\begin{enumerate}[leftmargin=*]
    \item \imp{气型污染}——大气污染物沉降
    \item \imp{水型污染}——工业废水和生活污水通过污水灌溉污染
    \item \imp{固体废弃物型污染}——垃圾、工业废渣等
\end{enumerate}

\subsection{土壤自净}

\begin{defbox}
\textbf{土壤自净作用（Soil Self-Purification）}：受污染土壤通过物理、化学和生物学作用，使病原体死灭，有害物质转化为无害，土壤逐渐恢复到污染前状态的过程。
\end{defbox}

\textbf{含氮有机物的无机化过程：}蛋白质等 → 氨基酸 → NH$_3$（\imp{氨化作用}）→ NO$_2^-$（亚硝化作用）→ NO$_3^-$（\imp{硝化作用}）

\subsection{重金属在土壤中的转化}

\begin{itemize}[leftmargin=*]
    \item 土壤胶体、腐殖质的吸附和螯合作用
    \item \imp{土壤pH}：pH越低，重金属溶解度越高，易被植物吸收
    \item 土壤氧化还原状态的影响
\end{itemize}

\section{土壤污染与健康}

\subsection{重金属污染}

\begin{keybox}
\textbf{镉（Cd）污染——痛痛病（Itai-itai Disease）：}
\begin{itemize}[leftmargin=*]
    \item 来源：矿山废水、冶炼、磷肥。水体镉污染主要来自工业废水；土壤镉污染主要是含镉废水灌溉农田所致
    \item IARC已将镉列为\imp{I类致癌物}。镉在人体中生物半减期长达\imp{10--25年}，可在体内不断累积。肾脏可蓄积镉吸收量的1/3，是镉中毒的主要靶器官
    \item 机制：镉与含羟基、氨基、巯基的蛋白质分子结合→抑制酶系统→影响肝肾功能；镉损伤肾小管→维生素D$_3$活性受影响→骨骼生长代谢受阻→钙磷代谢障碍→骨质疏松软化；镉干扰铁代谢→肠道铁吸收减低、破坏红细胞→贫血
    \item 表现：全身剧烈刺痛（腰、背、膝关节→遍及全身），骨骼严重畸形，骨脆易折，多发性病理骨折，肾功能损害。尿镉反映镉性肾损伤和体内镉负荷（以5 $\mu$mol/mol肌酐为慢性镉中毒诊断下限值），血镉反映近期接触量
    \item 流行病学：多为40岁以上妇女，妊娠、哺乳、营养不良等是发病诱因，常分布在镉矿丰富、雨量大、坡度陡、水土流失严重地区
\end{itemize}

\textbf{铊（Tl）污染：}神经系统和消化系统损害，特征性脱发

\textbf{铬（Cr）污染：}Cr(VI)毒性远大于Cr(III)；皮肤溃疡（铬疮）、呼吸道刺激、致癌性
\end{keybox}

\subsection{农药污染与POPs}

\begin{keybox}
\textbf{持久性有机污染物（POPs）的四大特性：}
\begin{enumerate}[leftmargin=*]
    \item \imp{持久性}——在环境中难以降解
    \item \imp{蓄积性}——易在生物体内脂肪组织中蓄积
    \item \imp{迁移性}——可通过大气、水流远距离迁移（"全球蒸馏效应"）
    \item \imp{高毒性}——具有致癌、致畸、致突变"三致"效应
\end{enumerate}

\textbf{健康危害：}急慢性毒性；"三致"作用；干扰生殖内分泌功能；免疫毒性、神经发育毒性。
\end{keybox}

\subsection{生物性污染的危害}

\begin{enumerate}[leftmargin=*]
    \item \imp{肠道传染病和寄生虫病}——经口传播
    \item \imp{钩端螺旋体病和炭疽病}——接触污染土壤或水
    \item \imp{破伤风和肉毒中毒}——伤口接触含芽孢的土壤
\end{enumerate}

\section{土壤污染的预防对策}

\begin{itemize}[leftmargin=*]
    \item \imp{土壤污染防治法律法规}
    \item \imp{土壤环境质量标准}
    \item \imp{土壤卫生防护}：粪便无害化处理、垃圾无害化处理（卫生填埋、焚烧、堆肥）、危险工业废渣处理、污水灌溉卫生防护
    \item \imp{污染土壤修复}：物理修复、化学修复、生物修复
\end{itemize}

\section{复习题}

\begin{enumerate}[leftmargin=*, itemsep=8pt]
    \item \textbf{【名词解释】}土壤环境背景值；环境容量；腐殖质
    \item \textbf{【名词解释】}土壤污染；土壤自净；POPs
    \item \textbf{【简答题】}简述土壤环境背景值的实践意义。
    \item \textbf{【简答题】}简述土壤污染的基本特点。
    \item \textbf{【简答题】}简述土壤自净的机制（以含氮有机物无机化为例）。
    \item \textbf{【简答题】}试述镉污染与痛痛病的关系（来源、机制、表现）。
    \item \textbf{【简答题】}POPs有哪些特性？对健康有何危害？
    \item \textbf{【简答题】}简述土壤生物性污染引起的主要疾病及危害途径。
    \item \textbf{【简答题】}简述土壤的卫生防护原则。
    \item \textbf{【非标准答案题】}结合痛痛病和水俣病的案例，谈谈食物链传递在环境污染健康危害中的关键作用。
\end{enumerate}

\subsection*{参考答案要点}

\begin{enumerate}[leftmargin=*, itemsep=5pt]
    \item \textbf{背景值}：未受污染天然土壤中各元素含量。\textbf{环境容量}：不超过卫生标准前提下污染物能容纳的最大负荷量。\textbf{腐殖质}：动植物残体经分解转化形成的有机物质。

    \item \textbf{土壤污染}：人类活动排放有害物质进入土壤危害人畜健康。\textbf{土壤自净}：受污染土壤通过物化生作用恢复的过程。\textbf{POPs}：具持久性、蓄积性、迁移性和高毒性的有机污染物。

    \item 四项意义：评价污染程度的参照值；制定卫生标准的依据；评价对健康影响的依据；土地开发和地方病防治的科学依据。

    \item 四大特点：隐蔽性、累积性与地域性、不可逆转性、治理长期性。

    \item 物理+化学+生物净化。含氮有机物无机化：蛋白质→氨基酸→NH$_3$（氨化）→NO$_2^-$（亚硝化）→NO$_3^-$（硝化）。

    \item 来源：矿山废水、冶炼、磷肥。机制：Cd蓄积→肾小管损伤→钙磷障碍→骨软化。表现：全身剧烈疼痛、骨折、肾功能衰竭。

    \item 四特性：持久性、蓄积性、迁移性、高毒性。危害：急慢性毒性、三致效应、干扰内分泌、免疫和神经毒性。

    \item 肠道传染病和寄生虫病（经口）、钩端螺旋体病和炭疽病（接触）、破伤风和肉毒中毒（伤口感染）。

    \item 法律标准制定、土壤卫生防护（粪便垃圾无害化等）、发展生态农业、污染土壤修复。

    \item 围绕食物链的生物放大效应、从环境到人体的传递路径、对风险预防和治理的启示展开。
\end{enumerate}

% ----- PPT参考题 -----
\subsection*{参考题（PPT课件补充）}

\subsubsection*{一、思考题}

\begin{enumerate}[leftmargin=*, itemsep=8pt]
    \item \textbf{【思考题】}某一地区农用土壤样品检测，11项指标（8种重金属+2种农药+苯并[a]芘）均低于风险筛选值，是否表明该土壤一定无污染？请说明理由。
\end{enumerate}

\subsubsection*{二、判断题（判断下列说法是否正确，错误的请说明理由）}

\begin{enumerate}[leftmargin=*, itemsep=8pt]
    \item 土壤污染引起破伤风的主要危害途径是土壤-人。（~~~~）
    \item 土壤中元素的本底值是指任何一处土壤中各种元素的含量。（~~~~）
    \item 土壤的pH较高，呈碱性时有利于农作物吸收重金属元素。（~~~~）
    \item 土壤污染的基本特点包括隐蔽性、累积性、迁移性、长期性。（~~~~）
    \item 持久性有机污染物特性包括持久性、迁移性、生物蓄积性、高毒性。（~~~~）
    \item 镉污染所致痛痛病属于土壤水型污染。（~~~~）
    \item 土壤污染的来源包括地质环境中区域性差异导致土壤中某些元素过高。（~~~~）
    \item 土壤颗粒越小，土壤容水性越大，透气性越差。（~~~~）
    \item 土壤气型污染的特点之一是污染物集中在土壤深层。（~~~~）
    \item 镉、汞、砷、铅、铬含量高于风险管制值的土地，原则上禁止种植食用农产品。（~~~~）
\end{enumerate}

\subsubsection*{参考答案}

\textbf{思考题：}
\begin{enumerate}[leftmargin=*, itemsep=5pt]
    \item 不一定。理由：①我国现行农用地土壤标准（GB 15618-2018）仅涵盖11项指标，但土壤污染物种类远多于此；②例如铊（Tl）已被加拿大、德国、美国等国家列为优先控制污染物，但我国尚未将其纳入土壤环境质量标准；③可能存在"假阴性"——未纳入标准的污染物即使超标也无法检出。因此，应关注标准外的潜在污染物。
\end{enumerate}

\textbf{判断题：}
\begin{enumerate}[leftmargin=*, itemsep=5pt]
    \item \textbf{×（错误）}。破伤风的主要危害途径是\textbf{土壤→伤口→人}，即含破伤风芽孢的土壤通过破损皮肤进入人体，而非直接"土壤-人"。
    \item \textbf{×（错误）}。本底值（背景值）是指\textbf{未受人为污染的天然土壤}中各种元素的含量，而非"任何一处土壤"。受污染地区的土壤元素含量不属于本底值。
    \item \textbf{×（错误）}。土壤pH越高（碱性越强），重金属溶解度越低，形成不溶性氢氧化物，反而\textbf{不利于}农作物吸收。酸性土壤中重金属活性更高，更易被作物吸收。
    \item \textbf{×（错误）}。土壤污染四大基本特点为：\textbf{隐蔽性、累积性与地域性、不可逆转性、治理长期性}。"迁移性"不属于基本特点（迁移性是POPs的特性）。
    \item \textbf{√（正确）}。POPs的四大特性为持久性、迁移性（可远距离传输）、生物蓄积性、高毒性。
    \item \textbf{√（正确）}。痛痛病由镉污染通过废水→土壤→稻米→人体的水型污染途径引起。
    \item \textbf{×（错误）}。地质环境中区域性差异导致土壤中某些元素过高属于\textbf{原生环境}问题（如生物地球化学性疾病），不属于人为活动造成的"土壤污染"。
    \item \textbf{√（正确）}。土壤颗粒越小，孔隙越小，容水性（持水性）越大，但透气性越差。
    \item \textbf{×（错误）}。气型污染主要由大气污染物沉降引起，污染物主要集中在土壤\textbf{表层}，而非深层。
    \item \textbf{√（正确）}。根据《农用地土壤污染风险管控标准》，重金属含量高于风险管制值的土地，原则上应采取禁止种植食用农产品等严格管控措施。
\end{enumerate}

\newpage

% =====================================================================
%  CHAPTER 8: 环境相关疾病
% =====================================================================
\newpage